\documentclass[12pt, a4paper]{article}
\usepackage[utf8]{inputenc}
\usepackage{geometry}
\usepackage{parskip}
\usepackage{amsmath}
\usepackage{microtype}

\geometry{top=2.5cm, bottom=2.5cm, left=2.5cm, right=2.5cm}

\title{Class 12 Chemistry NCERT Solutions}
\author{}
\date{}

\begin{document}

\maketitle

\section*{Chapter 1: Solutions}

\subsection*{Problem 1.1}
Calculate the mass percentage of benzene ($C_6H_6$) and carbon tetrachloride ($CCl_4$) if 22 g of benzene is dissolved in 122 g of carbon tetrachloride.

\textbf{Solution:}
Mass of solution = Mass of $C_6H_6$ + Mass of $CCl_4$
Mass of solution = $22\text{ g} + 122\text{ g} = 144\text{ g}$

Mass percentage of $C_6H_6 = \frac{\text{Mass of } C_6H_6}{\text{Total mass of solution}} \times 100$
Mass percentage of $C_6H_6 = \frac{22}{144} \times 100 = 15.28\%$

Mass percentage of $CCl_4 = \frac{\text{Mass of } CCl_4}{\text{Total mass of solution}} \times 100$
Mass percentage of $CCl_4 = \frac{122}{144} \times 100 = 84.72\%$

\subsection*{Problem 1.2}
State Raoult's Law.
\textbf{Answer:} Raoult's law states that for a solution of volatile liquids, the partial vapour pressure of each component of the solution is directly proportional to its mole fraction present in solution. For a binary solution of components A and B:
\[ P_A = P_A^0 \chi_A \]
\[ P_B = P_B^0 \chi_B \]

\section*{Chapter 2: Electrochemistry}

\subsection*{Problem 2.1}
Represent the cell in which the following reaction takes place:
$Mg(s) + 2Ag^+(0.0001 M) \rightarrow Mg^{2+}(0.130 M) + 2Ag(s)$
Calculate its $E_{cell}$ if $E^0_{cell} = 3.17$ V.

\textbf{Solution:}
The cell representation is: $Mg(s) | Mg^{2+}(0.130 M) || Ag^+(0.0001 M) | Ag(s)$
Using Nernst Equation:
\[ E_{cell} = E^0_{cell} - \frac{0.0591}{n} \log \frac{[Mg^{2+}]}{[Ag^+]^2} \]
Here $n = 2$.
\[ E_{cell} = 3.17 - \frac{0.0591}{2} \log \frac{0.130}{(0.0001)^2} \]
\[ E_{cell} = 3.17 - 0.02955 \log(1.3 \times 10^7) \]
\[ E_{cell} = 3.17 - 0.02955(7.1139) \]
\[ E_{cell} = 3.17 - 0.21 = 2.96 \text{ V} \]

\section*{Chapter 3: Chemical Kinetics}

\subsection*{Problem 3.1}
The rate constant for a first order reaction is $60 \text{ s}^{-1}$. How much time will it take to reduce the initial concentration of the reactant to its $1/16^{th}$ value?

\textbf{Solution:}
For a first order reaction:
\[ t = \frac{2.303}{k} \log \frac{[R]_0}{[R]} \]
Given $k = 60 \text{ s}^{-1}$ and $[R] = \frac{[R]_0}{16}$.
\[ t = \frac{2.303}{60} \log \frac{[R]_0}{[R]_0/16} \]
\[ t = \frac{2.303}{60} \log 16 \]
\[ t = \frac{2.303}{60} \log 2^4 = \frac{2.303 \times 4 \times 0.3010}{60} \]
\[ t = \frac{2.773}{60} = 0.0462 \text{ s} = 4.62 \times 10^{-2} \text{ s} \]

\section*{Chapter 4: d- and f- Block Elements}

\subsection*{Problem 4.1}
Why do transition elements exhibit variable oxidation states?
\textbf{Answer:} Transition elements exhibit variable oxidation states because the energy difference between $(n-1)d$ and $ns$ orbitals is very small. As a result, electrons from both energy levels can participate in bond formation.

\subsection*{Problem 4.2}
Calculate the 'spin-only' magnetic moment of $M^{2+}(aq)$ ion ($Z = 27$).
\textbf{Answer:} $Z = 27$ is Cobalt (Co). Electronic configuration of $Co = [Ar] 3d^7 4s^2$.
Electronic configuration of $Co^{2+} = [Ar] 3d^7$.
Number of unpaired electrons ($n$) in $3d^7$ ($t_{2g}^5 e_g^2$):
$d^7$ has 3 unpaired electrons.
Magnetic moment $\mu = \sqrt{n(n+2)}$ BM
$\mu = \sqrt{3(3+2)} = \sqrt{15} = 3.87$ BM.

\end{document}
